\subsection{诱导测度的性质}

命题 3. --- 设 X 是 T 的一个局部紧子空间，\(\lambda\) 是 X 上的一个复测度。下列性质等价：

\begin{enumerate}
\def\labelenumi{\alph{enumi})}
\item
  典范单射 \(i:\mathrm{X}\to \mathrm{T}\) 是 \(\lambda\) -逆紧的；
\item
  对 T 的每个紧子集 K，\(\mathbf{K} \cap \mathbf{X}\) 本质 \(\lambda\) -可积；
\item
  每个点 \(t \in \mathbf{T}\) 都有一个邻域 \(\mathbf{V}\) ，使得 \(\mathbf{V} \cap \mathbf{X}\) 本质 \(\lambda\) -可积；
\item
  存在 T 上的一个测度 \(\theta\) ，使得 \(\theta_{\mathrm{X}} = \lambda\) 。
\end{enumerate}

若这些等价条件满足，则沿用 d) 中的记号，有

\[
\left(i (\lambda)\right) _ {\mathrm{X}} = \lambda \quad \text{且} \quad i (\lambda) = i (\theta_ {\mathrm{X}}) = \varphi_ {\mathrm{X}} \cdot \theta .\tag{7}
\]

由于单射 \(i\) 是连续的，性质 a)、b)、c) 的等价性由 §6 的命题 1 及其后的注记应用于正测度 \(|\lambda|\) 而得。若 \(\lambda\) 在 X 上由 T 上的一个测度 \(\theta\) 诱导，则 \(|\lambda| = |\theta|_{\mathrm{X}}\) （公式 (6)），于是对 T 的每个紧子集 K 有 \(|\lambda| (\mathrm{K} \cap \mathrm{X}) = |\theta| (\mathrm{K} \cap \mathrm{X}) \leqslant |\theta| (\mathrm{K}) < +\infty\) （命题 1），故 d) 蕴涵 b)。设 a) 满足，让我们证明 \(\left(i(\lambda)\right)_{\mathrm{X}} = \lambda\) ，由此将得 d)。设 \(g\) 是 \(\mathcal{H}(\mathrm{X}; \mathbf{C})\) 的一个元素；用 \(g'\) 表示 \(g\) 到 T 上的零延拓，则依诱导测度的定义，再由 §6 第 4 号的命题 7，有

\[
\int g d (i (\lambda)) _ {X} = \int g ^ {\prime} d (i (\lambda)) = \int (g ^ {\prime} \circ i) d \lambda = \int g d \lambda .
\]

这就完成了四个性质等价性的证明。若 \(\lambda = \theta_{\mathrm{X}}\) 且 \(g\in \mathcal{K}(\mathrm{T};\mathbf{C})\) ，则

\[
\int g d \bigl (i (\theta_ {\mathrm{X}}) \bigr) = \int (g \circ i) d (\theta_ {\mathrm{X}}) = \int g \varphi_ {\mathrm{X}} d \theta ,
\]

因为 \(g\varphi_{\mathbf{X}}\) 是 \(g \circ i\) 到 \(\mathbf{T}\) 上的零延拓。这就证明了 (7) 的第二个公式。

推论 1. --- 若 X 是闭的，则 X 上的每个复测度 \(\lambda\) 都由 T 上的一个测度所诱导。

因为若 K 是 T 中的紧集，则 \(K \cap X\) 是紧的，从而是 \(\lambda\) -可积的。

推论 2. --- 设 \(\theta\) 是 T 上的一个复测度，\(\pi\) 是 T 到局部紧空间 Y 中的一个 \(\theta\) -逆紧映射，\(\pi_{X}\) 是它在 X 上的限制。则 \(\pi_{X}\) 是 \(\theta_{X}\) -逆紧的，且 \(\pi_{X}(\theta_{X}) = \pi(\varphi_{X} \cdot \theta)\) 。

因为 \(\pi_{X} = \pi \circ i\) ，其中 i 是典范单射 \(X \to T\) 。当 \(\theta\) 为正测度时，推论可由命题 3 及像测度的传递性（ \(\S6\) , No.~3, 命题 4）得出。非正复测度的情形随后由线性得出。

命题 4. --- 设 X 和 Y 是 T 的两个局部紧子空间，满足 \(Y \subset X\) 。若 \(\theta\) 是 T 上的一个复测度，则测度 \((\theta_{\mathrm{X}})_{\mathrm{Y}}\) （由 \(\theta_{X}\) 在 Y 上诱导）等于 \(\theta_{Y}\) （“诱导测度的传递性”）。

只须注意：若 g 是 \(\mathcal{K}(\mathrm{Y};\mathbf{C})\) 的一个元素，则 g 到 T 上的零延拓可以由把 g 到 X 上的零延拓再零延拓到 T 而得到；或者，利用概说中的等同，即 \(\varphi_{\mathrm{Y}}\cdot\theta=\varphi_{\mathrm{Y}}(\varphi_{\mathrm{X}}\cdot\theta)\) （§5, No.~4, 命题 8）。

命题 5. --- 设 \((\lambda_{\alpha})_{\alpha\in A}\) 是 T 上的正测度组成的递增有向族，具有上确界 \(\lambda\) ，并设 X 是 T 的一个局部紧子空间。则诱导测度族 \(\lambda_{\alpha}|X\) 在 \(\mathcal{M}(X)\) 中有上界，且

\[
\sup _ {\alpha \in \mathrm{A}} (\lambda_ {\alpha} | \mathrm{X}) = \lambda | \mathrm{X}.\tag{8}
\]

考虑到概说中的等同，本命题是 §5, No.~4 命题 5 的特殊情形。

推论. --- 设 \((\mu_{i})_{i\in\mathbb{I}}\) 是 T 上的正测度组成的可和族，其和为 \(\mu\) 。则诱导测度族 \(\mu_{i}|X\) 是可和的，且

\[
\sum_ {i \in I} (\mu_ {i} | X) = \mu | X.\tag{9}
\]

命题 6. --- 设 \(\Lambda: t \mapsto \lambda_t\) 是一个 \(\mu\) -适宜映射，把 T 映入 \(\mathcal{M}_+(\mathrm{X})\) ，其中 X 是在无穷远处可数的局部紧空间，并设 Y 是 X 的一个局部紧子空间。置 \(\int \lambda_t d\mu(t) = \nu\) 。则映射 \(t \mapsto \lambda_t|\mathrm{Y}\) （T 到 \(\mathcal{M}_+(\mathrm{Y})\) 中的映射）是 \(\mu\) -适宜的，且

\[
\int (\lambda_ {t} | \mathrm{Y}) d \mu (t) = \nu | \mathrm{Y}.\tag{10}
\]

考虑到概说中的等同，本命题是 §5, No.~4 命题 7 的特殊情形。

